\clearpage
\section*{Пятёрка 1. Вопросы 1–5}

% --- Вопрос 1 ---
\examquestion{1}{Метод интегрирования по частям в неопределенном интеграле}
{Проведите вывод формулы интегрирования по частям, опираясь на дифференциал произведения двух функций. Перечислите и классифицируйте три основных типа интегралов, берущихся данным методом (с примерами шаблонов для каждого типа).}

\paragraph{Теоретическое обоснование и вывод формулы.}
Пусть на некотором промежутке заданы две функции $u(x)$ и $v(x)$, непрерывно дифференцируемые на этом промежутке. По правилу дифференцирования произведения, дифференциал функции $u(x)v(x)$ имеет вид:
$$ d(uv) = u \, dv + v \, du $$
В силу операции интегрирования, которая является обратной к дифференцированию, непрерывные функции в правой части равенства имеют первообразные. Интегрируя левую и правой части, в силу свойства линейности неопределенного интеграла получаем:
$$ \int d(uv) = \int (u \, dv + v \, du) = \int u \, dv + \int v \, du $$
По свойству неопределенного интеграла от дифференциала функции, левая часть равна $uv$ (с точностью до константы). Таким образом:
$$ uv = \int u \, dv + \int v \, du $$
Выражая интеграл $\int u \, dv$, получаем классическую формулу интегрирования по частям:
$$ \int u \, dv = uv - \int v \, du $$
Смысл данной формулы заключается в том, что вычисление интеграла $\int u \, dv$ сводится к вычислению другого интеграла $\int v \, du$, который при правильном выборе функций $u$ и $v$ оказывается существенно проще исходного.

\paragraph{Классификация трех основных типов интегралов:}
\begin{enumerate}
    \item \textbf{Тип I (Понижение степени многочлена):} Интегралы вида $\int P_n(x) e^{ax} dx$, $\int P_n(x) \sin(bx) dx$, $\int P_n(x) \cos(bx) dx$, где $P_n(x)$ — многочлен степени $n$. В данном случае за $u$ целесообразно выбирать многочлен $u = P_n(x)$, так как при его дифференцировании степень уменьшается ($du = P_n'(x)dx$), а в качестве $dv$ берется оставшаяся тригонометрическая или показательная часть, легко интегрируемая методом табличных значений. Шаг повторяется $n$ раз до полного исчезновения многочлена.
    
    \item \textbf{Тип II (Переход к алгебраическим дробям):} Интегралы вида $\int P_n(x) \ln x \, dx$, $\int P_n(x) \arcsin x \, dx$, $\int P_n(x) \arctan x \, dx$. Здесь за $u$ выбирается трансцендентная функция ($u = \ln x$ или $u = \arctan x$), поскольку операция дифференцирования превращает её в алгебраическое выражение. В качестве $dv$ берется многочлен с дифференциалом $dv = P_n(x)dx$, интегрирование которого дает функцию $v$, являющуюся многочленом степени $n+1$. После применения формулы интегрирования по частям новый интеграл представляет собой интегрирование рациональной дроби.
    
    \item \textbf{Тип III (Циклические интегралы):} Интегралы вида $\int e^{ax} \sin(bx) dx$ и $\int e^{ax} \cos(bx) dx$. Особенность данного типа заключается в том, что ни показательная, ни тригонометрическая функции не упрощаются при дифференцировании или интегрировании. Применяя формулу по частям дважды (выбирая на каждом шаге за $u$ функции одного класса — либо всегда экспоненту, либо всегда тригонометрию), мы возвращаемся к исходному интегралу, формируя линейное уравнение относительно искомого интеграла.
\end{enumerate}

\paragraph{Бонус-пример на каждый тип:}
\begin{itemize}
    \item \textbf{Пример для Типа I:} Вычислим интеграл $\int x e^{2x} dx$. Пусть $u = x \implies du = dx$, а $dv = e^{2x} dx \implies v = \frac{1}{2}e^{2x}$. Тогда:
    $$ \int x e^{2x} dx = x \cdot \left(\frac{1}{2}e^{2x}\right) - \int \frac{1}{2}e^{2x} dx = \frac{x}{2}e^{2x} - \frac{1}{4}e^{2x} + C = \frac{e^{2x}}{4}(2x - 1) + C. $$
    
    \item \textbf{Пример для Типа II:} Вычислим интеграл $\int \ln x \, dx$. Пусть $u = \ln x \implies du = \frac{1}{x}dx$, а $dv = dx \implies v = x$. Применяем формулу:
    $$ \int \ln x \, dx = x \ln x - \int x \cdot \frac{1}{x} dx = x \ln x - \int dx = x \ln x - x + C. $$
    
    \item \textbf{Пример для Типа III:} Найдем интеграл $I = \int e^x \cos x \, dx$. Пусть $u = e^x \implies du = e^x dx$, а $dv = \cos x \, dx \implies v = \sin x$. Получаем:
    $$ I = e^x \sin x - \int e^x \sin x \, dx $$
    Применим метод по частям к полученному интегралу повторно, удерживая выбор $u$ за экспонентой: $u = e^x \implies du = e^x dx$, $dv = \sin x \, dx \implies v = -\cos x$.
    $$ I = e^x \sin x - \left( -e^x \cos x - \int (-\cos x) e^x dx \right) = e^x \sin x + e^x \cos x - \int e^x \cos x \, dx $$
    Замечаем, что последний интеграл равен $I$. Получаем алгебраическое уравнение:
    $$ I = e^x(\sin x + \cos x) - I \implies 2I = e^x(\sin x + \cos x) \implies I = \frac{e^x(\sin x + \cos x)}{2} + C. $$
\end{itemize}







% --- Вопрос 2 ---
\examquestion{2}{Метод замены переменной (подстановки) в неопределенном интеграле}
{Сформулируйте и докажите теорему о замене переменной. Покажите на конкретном примере, как этот метод применяется для раскрытия циклических тригонометрических интегралов.}

\begin{theorem}[О замене переменной]
Пусть функция $x = \varphi(t)$ определена и дифференцируема на промежутке $T$, а её множество значений представляет собой промежуток $X$. Пусть функция $f(x)$ определена на промежутке $X$ и имеет на нем первообразную $F(x)$, то есть $F'(x) = f(x)$. Тогда на промежутке $T$ функция $f(\varphi(t))\varphi'(t)$ имеет первообразную $F(\varphi(t))$, и справедлива формула:
$$ \int f(\varphi(t)) \varphi'(t) dt = \left. \int f(x) dx \right|_{x = \varphi(t)} = F(\varphi(t)) + C $$
\end{theorem}

\begin{proof}
Для доказательства данного тождества достаточно показать, что функция $F(\varphi(t))$ является первообразной для подынтегральной функции в левой части равенства. Согласно определению, продифференцируем сложную функцию $F(\varphi(t))$ по переменной $t$, используя правило дифференцирования сложной функции:
$$ \frac{d}{dt} \Big( F(\varphi(t)) \Big) = F'(\varphi(t)) \cdot \varphi'(t) $$
Так как по условию теоремы $F'(x) = f(x)$, то, подставляя аргумент $x = \varphi(t)$, имеем $F'(\varphi(t)) = f(\varphi(t))$. Отсюда получаем:
$$ \frac{d}{dt} \Big( F(\varphi(t)) \Big) = f(\varphi(t)) \varphi'(t) $$
Поскольку производная функции $F(\varphi(t))$ по переменной $t$ в точности совпадает с подынтегральной функцией, то по определению неопределенного интеграла множество всех первообразных имеет вид $F(\varphi(t)) + C$, что завершает доказательство теоремы.
\end{proof}

\paragraph{Пример (Раскрытие циклического тригонометрического интеграла).}
Рассмотрим интеграл от нечетной степени тригонометрической функции, порождающий циклический переход при неверном подходе, но эффективно раскрываемый подстановкой: $I = \int \sin^3 x \cos^2 x \, dx$.
Используем основное тригонометрическое тождество, чтобы выделить один дифференциал синуса: $\sin^3 x = \sin^2 x \cdot \sin x = (1 - \cos^2 x)\sin x$. Тогда:
$$ I = \int (1 - \cos^2 x)\cos^2 x \cdot \sin x \, dx $$
Сделаем замену переменной $t = \cos x$. Находим дифференциал: $dt = -\sin x \, dx \implies \sin x \, dx = -dt$. Пересчитаем интеграл относительно переменной $t$:
$$ I = \int (1 - t^2)t^2 (-dt) = \int (t^4 - t^2) dt = \frac{t^5}{5} - \frac{t^3}{3} + C $$
Возвращаясь к исходной переменной $x$, получаем окончательный ответ:
$$ I = \frac{\cos^5 x}{5} - \frac{\cos^3 x}{3} + C. $$

\paragraph{Бонус-пример (Универсальная тригонометрическая подстановка):}
Вычислим $I = \int \frac{dx}{2 + \cos x}$. Сделаем стандартную подстановку $t = \tan \frac{x}{2}$. Тогда выразим $dx = \frac{2dt}{1+t^2}$ и $\cos x = \frac{1-t^2}{1+t^2}$. Подставляем в интеграл:
$$ I = \int \frac{\frac{2dt}{1+t^2}}{2 + \frac{1-t^2}{1+t^2}} = \int \frac{2dt}{2(1+t^2) + (1-t^2)} = \int \frac{2dt}{3 + t^2} = \frac{2}{\sqrt{3}} \arctan\left(\frac{t}{\sqrt{3}}\right) + C $$
Выполняя обратную замену, находим: $I = \frac{2}{\sqrt{3}} \arctan\left(\frac{\tan(x/2)}{\sqrt{3}}\right) + C$.





% --- Вопрос 3 ---
\examquestion{3}{Complex method v veschestvennom analize}
{Опираясь на свойства комплексной экспоненты и формулу Эйлера, проведите вывод формул для интегралов вида $I_{1}=\int e^{ax}\cos(bx)dx$ и $I_{2}=\int e^{ax}\sin(bx)dx$, где $a,b\in\mathbb{R}$. Докажите корректность перехода к комплексному интегралу $\int e^{(a+ib)x}dx$ и последующего разделения вещественной и мнимой частей.}

\paragraph{Суть метода и итоговые формулы.}
При классическом подходе интегралы вида $I_1 = \int e^{ax}\cos(bx)dx$ и $I_2 = \int e^{ax}\sin(bx)dx$ вычисляются двукратным интегрированием по частям, что приводит к громоздким циклическим уравнениям. 

Комплексный метод позволяет найти оба интеграла одновременно и чисто алгебраическим путём. Для этого вводится вспомогательный комплексный интеграл $J = \int e^{(a+ib)x} dx$. Согласно формуле Эйлера, его вещественная часть ($\text{Re}\,J$) в точности равна интегралу $I_1$, а мнимая часть ($\text{Im}\,J$) — интегралу $I_2$.

В результате интегрирования комплексной экспоненты получаются следующие **итоговые формулы**:
$$ \int e^{ax}\cos(bx)dx = \frac{e^{ax}(a\cos bx + b\sin bx)}{a^2+b^2} + C $$
$$ \int e^{ax}\sin(bx)dx = \frac{e^{ax}(a\sin bx - b\cos bx)}{a^2+b^2} + C $$
Ниже приведено строгое обоснование этого перехода и пошаговый алгебраический вывод.

\paragraph{Доказательство корректности перехода.}
Пусть задана комплексная функция вещественного аргумента $w(x) = u(x) + i v(x)$, где $u(x), v(x)$ — вещественные функции, непрерывные на рассматриваемом промежутке. По определению, интеграл от такой функции равен:
$$ \int w(x) dx = \int (u(x) + i v(x)) dx = \int u(x) dx + i \int v(x) dx $$
Формула Эйлера задает связь комплексной экспоненты с тригонометрическими функциями: $e^{i\varphi} = \cos\varphi + i\sin\varphi$. Применим её к функции $e^{(a+ib)x}$:
$$ e^{(a+ib)x} = e^{ax} \cdot e^{ibx} = e^{ax}(\cos(bx) + i \sin(bx)) = e^{ax}\cos(bx) + i e^{ax}\sin(bx) $$
Интегрируя это выражение по определению, получаем:
$$ \int e^{(a+ib)x} dx = \int e^{ax}\cos(bx) dx + i \int e^{ax}\sin(bx) dx = I_1 + i I_2 $$
Поскольку операция выделения вещественной ($\text{Re}$) и мнимой ($\text{Im}$) частей коммутирует с интегрированием, мы можем утверждать:
$$ I_1 = \text{Re} \left( \int e^{(a+ib)x} dx \right), \quad I_2 = \text{Im} \left( \int e^{(a+ib)x} dx \right) $$
Так как закон дифференцирования $\frac{d}{dx} e^{kx} = k e^{kx}$ выполняется на комплексной плоскости для любого $k = a+ib \in \mathbb{C}$, операция вычисления первообразной от комплексной экспоненты полностью эквивалентна вещественному аналогу: $\int e^{kx} dx = \frac{1}{k} e^{kx} + C$.

\paragraph{Вывод формул.}
Вычислим комплексный интеграл в явном виде:
$$ \int e^{(a+ib)x} dx = \frac{e^{(a+ib)x}}{a+ib} = \frac{e^{ax}(\cos bx + i \sin bx)}{a+ib} $$
Для разделения вещественной и мнимой частей избавимся от комплексности в знаменателе, умножив числитель и знаменатель на сопряженное число $(a-ib)$:
$$ \frac{e^{ax}(\cos bx + i \sin bx)(a-ib)}{(a+ib)(a-ib)} = \frac{e^{ax}}{a^2+b^2} \Big( a\cos bx - i b\cos bx + i a\sin bx - i^2 b\sin bx \Big) $$
Учитывая, что $i^2 = -1$, сгруппируем слагаемые на вещественные и мнимые компоненты:
$$ \int e^{(a+ib)x} dx = \frac{e^{ax}}{a^2+b^2} \Big( (a\cos bx + b\sin bx) + i(a\sin bx - b\cos bx) \Big) $$
Выделяя вещественную и мнимую части, получаем искомые формулы для интегралов:
$$ I_1 = \int e^{ax}\cos(bx)dx = \frac{e^{ax}(a\cos bx + b\sin bx)}{a^2+b^2} + C $$
$$ I_2 = \int e^{ax}\sin(bx)dx = \frac{e^{ax}(a\sin bx - b\cos bx)}{a^2+b^2} + C $$

\paragraph{Бонус-пример:} Вычислим интеграл $\int e^{3x} \sin(2x) dx$. 
В данном случае параметры $a = 3$, $b = 2$. Прямым подстановлением в выведенную формулу для $I_2$ получаем:
$$ \int e^{3x} \sin(2x) dx = \frac{e^{3x}(3\sin 2x - 2\cos 2x)}{3^2+2^2} + C = \frac{e^{3x}(3\sin 2x - 2\cos 2x)}{13} + C. $$








% --- Вопрос 4 ---
\examquestion{4}{Методы нахождения неопределенных коэффициентов}
{Проведите сравнительный анализ метода приравнивания коэффициентов при одинаковых степенях и метода частных значений. Дайте обоснование методу вычетов Хевисайда для случая простых линейных множителей в знаменателе.}

\paragraph{Правила формирования структуры числителей элементарных дробей.}
При разложении правильной рациональной дроби $\frac{P(x)}{Q(x)}$ вид многочлена в числителе каждой элементарной дроби однозначно определяется структурой сомножителей в знаменателе:
\begin{itemize}
    \item \textbf{Принцип формирования:} Степень многочлена в числителе должна быть строго на единицу меньше степени базового многочлена, находящегося внутри скобки знаменателя (без учета внешней кратности этой скобки).
    \item \textbf{Линейный сомножитель} $(x - a)^k$: Внутреннее выражение $x - a$ является многочленом первой степени ($x^1$). Следовательно, в числителе формируется многочлен нулевой степени — константа ($A$). Количество дробей в группе равно кратности $k$:
    $$ \frac{A_1}{x-a} + \frac{A_2}{(x-a)^2} + \dots + \frac{A_k}{(x-a)^k} $$
    \item \textbf{Квадратичный сомножитель} $(x^2 + px + q)^m$ с отрицательным дискриминантом ($D < 0$): Внутреннее выражение является многочленом второй степени ($x^2$). Следовательно, в числителе формируется линейный двухчлен — многочлен первой степени ($Bx + C$). Количество дробей равно кратности $m$:
    $$ \frac{B_1 x + C_1}{x^2 + px + q} + \frac{B_2 x + C_2}{(x^2 + px + q)^2} + \dots + \frac{B_m x + C_m}{(x^2 + px + q)^m} $$
\end{itemize}

\paragraph{Теоретическое обоснование разложения.}
Любая правильная рациональная дробь $\frac{P(x)}{Q(x)}$ над полем действительных чисел допускает единственное разложение на сумму элементарных дробей указанных типов. 
\begin{itemize}
    \item Каждому действительному корню $x = a$ кратности $k$ соответствует группа из $k$ дробей вида: $\sum_{i=1}^{k} \frac{A_i}{(x-a)^i}$, где $A_i \in \mathbb{R}$.
    \item Каждой паре комплексно-сопряженных корней кратности $m$ соответствует группа из $m$ дробей вида: $\sum_{j=1}^{m} \frac{B_j x + C_j}{(x^2+px+q)^j}$, где $B_j, C_j \in \mathbb{R}$.
\end{itemize}

\paragraph{Сравнительный анализ методов определения коэффициентов.}
После приведения всех элементарных дробей к общему знаменателю $Q(x)$ фиксируется тождественное равенство исходного числителя $P(x)$ и полученного многочлена с неопределенными коэффициентами $M(x)$: $P(x) \equiv M(x)$. Для нахождения неизвестных применяются два метода:

\begin{enumerate}
    \item \textbf{Метод частных значений:}
    \\ Основан на свойстве тождественности: равенство справедливо при любых значениях переменной $x$. В уравнение подставляются действительные корни знаменателя $Q(x)$, что приводит к занулению большинства слагаемых и мгновенному нахождению части коэффициентов.
    \\ \textit{Преимущества:} Высокая скорость вычислений. При наличии только простых действительных корней исключает необходимость решения систем уравнений.
    \\ \textit{Недостатки:} При наличии комплексных или кратных корней метод не позволяет напрямую определить все неизвестные. Требуется подстановка искусственных значений (например, $x=0$, $x=1$) и комбинирование с другими подходами.
    
    \item \textbf{Метод приравнивания коэффициентов:}
    \\ Основан на теореме о равенстве многочленов: два многочлена тождественно равны тогда и только тогда, когда равны их коэффициенты при одинаковых степенях переменной $x^k$. Проводится полное раскрытие скобок в $M(x)$, группировка по степеням $x$ и составление системы линейных алгебраических уравнений (СЛАУ).
    \\ \textit{Преимущества:} Абсолютная универсальность. Метод применим для любых типов и кратностей корней знаменателя.
    \\ \textit{Недостатки:} Высокая вычислительная трудоемкость. При больших степенях знаменателя приводит к громоздким СЛАУ, требующим применения метода Гаусса.
\end{enumerate}

\paragraph{Обоснование метода вычетов Хевисайда.}
Пусть рациональная дробь имеет вид $\frac{P(x)}{Q(x)}$, где знаменатель содержит простой линейный корень $x = x_1$ кратности $1$. Тогда многочлен знаменателя представим в виде $Q(x) = (x - x_1)R(x)$, причем $R(x_1) \neq 0$. 

Выделим из общего разложения дробь с искомым коэффициентом $A_1$:
$$ \frac{P(x)}{(x-x_1)R(x)} = \frac{A_1}{x-x_1} + \sum \text{оставшиеся дроби} $$
Умножим обе части тождества на линейный множитель $(x - x_1)$:
$$ \frac{P(x)}{R(x)} = A_1 + (x - x_1) \cdot \sum \text{оставшиеся дроби} $$
Перейдем к пределу при $x \to x_1$. Поскольку сумма оставшихся дробей представляет собой правильную рациональную функцию, не имеющую особенностей в точке $x_1$, произведение второго слагаемого на $(x-x_1)$ стремится к нулю. Получаем явное выражение для коэффициента:
$$ A_1 = \lim_{x \to x_1} \frac{P(x)}{R(x)} = \frac{P(x_1)}{R(x_1)} $$
Применим правило дифференцирования произведения к знаменателю: $Q'(x) = R(x) + (x - x_1)R'(x)$. При подстановке значения $x = x_1$ второе слагаемое обращается в нуль: $Q'(x_1) = R(x_1)$. Таким образом, получаем классическую предельную формулу Хевисайда:
$$ A_1 = \frac{P(x_1)}{Q'(x_1)} $$

\paragraph{Пример практического применения (Метод Хевисайда):}
Разложим правильную дробь $\frac{2x + 3}{(x-2)(x+3)} = \frac{A}{x-2} + \frac{B}{x+3}$.
\begin{itemize}
    \item Нахождение коэффициента $A$ при корне $x = 2$ путем исключения сомножителя $(x-2)$ из исходного знаменателя:
    $$ A = \left. \frac{2x + 3}{x + 3} \right|_{x=2} = \frac{2(2) + 3}{2 + 3} = \frac{7}{5} $$
    \item Нахождение коэффициента $B$ при корне $x = -3$ путем исключения сомножителя $(x+3)$ из исходного знаменателя:
    $$ B = \left. \frac{2x + 3}{x - 2} \right|_{x=-3} = \frac{2(-3) + 3}{-3 - 2} = \frac{-3}{-5} = \frac{3}{5} $$
\end{itemize}
Результат разложения: $\frac{2x + 3}{(x-2)(x+3)} = \frac{7}{5(x-2)} + \frac{3}{5(x+3)}$.

% --- Вопрос 5 ---
\examquestion{5}{Интегрирование элементарных дробей IV типа}
{Выведите рекуррентную формулу для вычисления интеграла от элементарной дроби IV типа: $I_{n}=\int\frac{dx}{(x^{2}+a^{2})^{n}}$, $n\in\mathbb{N}$. Покажите этап интегрирования по частям и алгоритм понижения степени знаменателя.}

\paragraph{План вывода кратко:}
\begin{enumerate}
    \item \textbf{Подмена в числителе:} $1 = \frac{1}{a^2} \cdot \frac{(x^2+a^2) - x^2}{(x^2+a^2)^n} \implies$ разбиваем на $I_{n-1}$ и «хвост» с $x^2$.
    \item \textbf{Части для «хвоста»:} Интеграл $\int \frac{x \cdot x \, dx}{(x^2+a^2)^n}$ берем по частям:
    \begin{itemize}
        \item $u = x \implies du = dx$
        \item $dv = \frac{x \, dx}{(x^2+a^2)^n} \implies v = \frac{(x^2+a^2)^{-n+1}}{2(1-n)}$
    \end{itemize}
    \item \textbf{Сборка:} Подставляем всё в уравнение $I_n = \frac{1}{a^2}I_{n-1} - \frac{1}{a^2}I^*$ и приводим подобные при $I_{n-1}$.
\end{enumerate}

\begin{tcolorbox}[colback=white, colframe=black, boxrule=0.5pt, arc=0pt, borderline={0.5pt}{0pt}{black, dashed}]
    \centering
    $$ I_n = \frac{x}{2a^2(n-1)(x^2+a^2)^{n-1}} + \frac{2n-3}{2a^2(n-1)} I_{n-1} $$
\end{tcolorbox}
\paragraph{Математический вывод формулы.}
Рассмотрим интеграл элементарной дроби IV типа: $I_n = \int \frac{dx}{(x^2+a^2)^n}$. Реализуем первый этап алгоритма:
$$ I_n = \frac{1}{a^2} \int \frac{a^2}{(x^2+a^2)^n} dx = \frac{1}{a^2} \int \frac{(x^2+a^2) - x^2}{(x^2+a^2)^n} dx $$
Используя свойство линейности неопределенного интеграла, разобьем выражение на два отдельных слагаемых:
$$ I_n = \frac{1}{a^2} \int \frac{x^2+a^2}{(x^2+a^2)^n} dx - \frac{1}{a^2} \int \frac{x^2}{(x^2+a^2)^n} dx $$
После сокращения первой дроби получаем соотношение, связывающее $I_n$ и $I_{n-1}$:
$$ I_n = \frac{1}{a^2} I_{n-1} - \frac{1}{a^2} I^*, \quad \text{где } I^* = \int \frac{x \cdot x \, dx}{(x^2+a^2)^n} $$
Перейдем ко второму этапу — интегрированию по частям интеграла $I^*$. Пусть:
$$ u = x \implies du = dx $$
$$ dv = \frac{x \, dx}{(x^2+a^2)^n} = \frac{1}{2} (x^2+a^2)^{-n} d(x^2+a^2) \implies v = \frac{1}{2(1-n)(x^2+a^2)^{n-1}} $$
Применяя классическую формулу $\int u \, dv = uv - \int v \, du$, преобразуем интеграл $I^*$:
$$ I^* = \frac{x}{2(1-n)(x^2+a^2)^{n-1}} - \int \frac{dx}{2(1-n)(x^2+a^2)^{n-1}} $$
Вынесем константу за знак интеграла, сводя его к определению $I_{n-1}$:
$$ I^* = \frac{x}{2(1-n)(x^2+a^2)^{n-1}} - \frac{1}{2(1-n)} I_{n-1} $$
Перейдем к третьему этапу. Подставим полученное выражение для $I^*$ в исходное уравнение для $I_n$:
$$ I_n = \frac{1}{a^2} I_{n-1} - \frac{1}{a^2} \left[ \frac{x}{2(1-n)(x^2+a^2)^{n-1}} - \frac{1}{2(1-n)} I_{n-1} \right] $$
Для упрощения знаков преобразуем константу в знаменателе: $\frac{1}{2(1-n)} = -\frac{1}{2(n-1)}$. Раскроем скобки:
$$ I_n = \frac{x}{2a^2(n-1)(x^2+a^2)^{n-1}} + \frac{1}{a^2} I_{n-1} - \frac{1}{2a^2(n-1)} I_{n-1} $$
Сгруппируем слагаемые при интеграле $I_{n-1}$, приведя их коэффициенты к общему знаменателю $2a^2(n-1)$:
$$ \frac{1}{a^2} - \frac{1}{2a^2(n-1)} = \frac{2(n-1) - 1}{2a^2(n-1)} = \frac{2n - 3}{2a^2(n-1)} $$
В результате получаем окончательную рекуррентную формулу понижения степени для интегралов IV типа:
$$ I_n = \frac{x}{2a^2(n-1)(x^2+a^2)^{n-1}} + \frac{2n-3}{2a^2(n-1)} I_{n-1} $$

\paragraph{Пример практического применения:} Вычислим интеграл $I_2 = \int \frac{dx}{(x^2+4)^2}$.
В данном случае параметры имеют значения $a^2 = 4 \implies a = 2$, начальная степень $n = 2$. 
Базовый табличный интеграл первого порядка ($n=1$) равен:
$$ I_1 = \int \frac{dx}{x^2+4} = \frac{1}{2} \arctan\left(\frac{x}{2}\right) $$
Подставим параметры $n=2$ и $a^2=4$ в выведенную рекуррентную формулу:
$$ I_2 = \frac{x}{2 \cdot 4 \cdot (2-1)(x^2+4)^{2-1}} + \frac{2(2)-3}{2 \cdot 4 \cdot (2-1)} I_1 = \frac{x}{8(x^2+4)} + \frac{1}{8} I_1 $$
Подставляя явный вид интеграла $I_1$ и добавляя константу интегрирования, получаем финальный результат:
$$ I_2 = \frac{x}{8(x^2+4)} + \frac{1}{16} \arctan\left(\frac{x}{2}\right) + C $$